\documentclass[10pt,a4paper]{amsart}
\usepackage[margin=19mm]{geometry}
\usepackage{amsmath,amssymb,mathtools}
\usepackage[scr=rsfs,cal=euler]{mathalfa}
\usepackage{xfrac}
\usepackage[unicode,hidelinks,bookmarks=false]{hyperref}
\newcommand{\ie}{i.e.\ }
\newcommand{\cf}{cf.\ }
\newcommand{\resp}{respectively\ }
\newcommand{\Subsection}{\subsection}
\providecommand{\nobreakdash}{\nobreak\hbox{-}}
\setlength{\emergencystretch}{2em}
\multlinegap0pt
\usepackage{amssymb}
\usepackage{bm}
\usepackage{float}
\usepackage{enumerate}
\newtheorem{lemma}{Lemma}
\title{Remarks on the parametric Schr\"{o}dinger-Poisson system}
\author{Chen Huang, Sihua Liang, Lei Ma, Patrizia Pucci}
\date{Source: arXiv:2508.12732v4 (CC BY 4.0)}
\begin{document}
\maketitle

\noindent\emph{Source and licence.} Remarks on the parametric Schr\"{o}dinger-Poisson system. Version arXiv:2508.12732v4 (CC BY 4.0), distributed by arXiv under the Creative Commons Attribution 4.0 International licence, http://creativecommons.org/licenses/by/4.0/, as recorded in the arXiv licence metadata for this exact version and shown on the abstract page https://arxiv.org/abs/2508.12732v4. Full licence notice: \texttt{licenses/arxiv-2508.12732-CC-BY-4.0.txt}.\par\smallskip

\noindent\textit{Editorial scope.} Counted unit: the article's lemma lem4 (source0.tex lines 727--729). The article's complete original proof of it is delivered with it (source0.tex lines 730--875, one \texttt{proof} environment of 146 lines). No mathematics is written, repaired or re-derived: the statement and the proof are the article's own lines, in the article's own notation, and no retained line is substituted or re-flowed.  Retained uncounted as the source-local context the proof needs: lines 698--707.  Nothing was omitted from any retained span, so the package carries no omission record.  No retained line contains a citation, so the wrapper carries no bibliography and no bibliography entry.  No retained line contains a figure environment, an image-inclusion command or drawing code, so no image is delivered and the package declares no asset dependency.  Editorial formatting only, in addition: an AMS \texttt{amsart} preamble that restates the article's own declarations and macros used by the retained lines, namely the environments \texttt{lemma} on separate counters, together with the standard packages those lines need.  The retained environments print in this document as Lemma 1 in order of appearance, whereas the article prints the same items with the numbering of its own full document.  The complete native source of the article as distributed in the arXiv e-print for this exact version is bundled unchanged as \texttt{sources/183432/source0.tex}.
\begin{lemma}\label{lem3.5}
       For any $R> 1$, there exists a constant $C_{R} > 0$ independent of $\epsilon $, such that for all $\left ( \epsilon ,u \right )\in \left [ 0,1 \right ]\times H_{V}^{1}\left (  \mathbb{R}^{3}\right )$, one has
      \begin{equation*}
      \left | I_{\epsilon,R }\left ( u \right )-I_{0}\left ( u \right )\right |\leq \frac{1}{R}+C_{R}\epsilon ^{p-2}
      \end{equation*}
      and
      \begin{equation*}
      \left | \langle I_{\epsilon,R }'\left ( u \right ),u\rangle-\langle I_{0}'\left ( u \right ),u\rangle\right |\leq \frac{1}{R}+C_{R}\epsilon ^{p-2}.
        \end{equation*}
\end{lemma}

\begin{lemma}\label{lem4}
Let $\{u_n\}_{n=1}^\infty\subset H^1_V(\mathbb R^3)$ be a Cerami sequence for the functional $I_{\epsilon,R}$. For any $R>1$ fixed, there exists a constant $\epsilon(R)>0$ such that if $\epsilon\in(0,\epsilon(R))$, then $u_n\rightarrow u\text{ in }H^1_V(\mathbb R^3)$ with $u\in H^1_V(\mathbb R^3)$.
\end{lemma}

\begin{proof}
    Let $\{u_n\}$ be a Cerami sequence of $I_{\epsilon,R}$, that is, for some $c\in \mathbb{R}$,
    $$I_{\epsilon,R}(u_n)\to c~\text{and}~(1+\|u_n\|_{H^{1}_{V}(\mathbb R^3)})I_{\epsilon,R}'(u_n)\to0\ ~\text{as}~n\to\infty.$$

Firstly, we claim that
$$\rho_n:=\|u_n\|_{H^{1}_{V}(\mathbb R^3)}<\infty.$$
If it is not true, we may assume $\rho_n\to\infty$. Let $v_n=\frac{u_n}{\rho_n}$. Consequently, $\{v_n\}$ is bounded in $H^{1}_{V}(\mathbb{R}^3)$. Up to a subsequence, $H^{1}_{V}(\mathbb{R}^{3})$ is compactly embedded in $L^2(\mathbb{R}^3)$, we may assume that
\begin{equation*}
		\begin{split}
	v_n&\to v~\text{weakly in}~H^{1}_{V}(\mathbb{R}^3);\\
	v_n&\to v~\text{strongly in}~L^{2}(\mathbb{R}^3);\\
	v_n&\to v~\text{a.e. in}~ \mathbb{R}^3.
	\end{split}\end{equation*}
Consider
\begin{equation*}
		\begin{split}
		c+o_{n}(1)&=I_{\epsilon,R}(u_n)-\frac{1}{3}\langle I_{\epsilon,R}'(u_n),u_n\rangle\\
		&=\left(\frac{1}{2}-\frac{1}{3}\right)\int_{\mathbb R^3}\left(|\nabla u_n|^2+\tilde V(x)u_n^2\right)\text{d}x+\left(\frac{1}{3}-\frac{1}{4}\right)\int_{\mathbb R^3}\phi_{u_n}u_n^2\text{d}x\\
        &\quad\quad-\frac{1}{\epsilon^2}\eta _{R}\left (\|u_n\|^{2}_{H^1_{V}(\mathbb{R}^3)}\right )\int_{\mathbb{R}^{3}}F\left (\epsilon \eta _{R} \left ( \left | u_n\right |\right )u_n\right )\text{d}x-\left(\frac{1}{2}-\frac{1}{3}\right)\int_{\mathbb{R}^3}m u^{2}_n\text{d}x\\
		&\quad\quad-\frac{2}{3\epsilon^{2}}\eta _{R}'\left ( \|u_n\|^{2}_{H^1_{V}(\mathbb{R}^3)} \right )\int_{\mathbb{R}^{3}}(|\nabla u_n|^2+\tilde V(x)u_n^2) \text{d}x\int_{\mathbb{R}^{3}}F\left (\epsilon\eta _{R}\left ( \left | u_n\right | \right ) u_n\right )\text{d}x\\
        &\quad\quad-\frac{1}{3\epsilon}\eta _{R}\left ( \|u_n\|^{2}_{H^1_{V}(\mathbb{R}^3)} \right )\int_{\mathbb{R}^{3}}\left(f\left (\epsilon \eta _{R}\left ( \left | u_n\right | \right ) u_n\right )\left (\eta _{R}\left ( \left | u_n\right | \right ) u_n +\eta _{R}'\left ( \left | u_n\right |\right )u_n^2 \right )\right)\text{d}x\\
		&\geq\left(\frac{1}{2}-\frac{1}{3}\right)\int_{\mathbb R^3}\left(|\nabla u_n|^2+\tilde V(x)u_n^2\right)\text{d}x+\left(\frac{1}{3}-\frac{1}{4}\right)\int_{\mathbb R^3}\phi_{u_n}u_n^2\text{d}x\\
		&\quad\quad-\frac{4}{3}\left(\frac{1}{R}+C_{R}\epsilon ^{p-2}\right)-\left(\frac{1}{2}-\frac{1}{3}\right)\int_{\mathbb{R}^3}m u^{2}_n\text{d}x\\
		&=\frac{1}{6}\int_{\mathbb R^3}\left(|\nabla u_n|^2+\tilde V(x)u_n^2\right)\text{d}x-\frac{1}{6}\int_{\mathbb{R}^3}m u^{2}_n\text{d}x-\frac{4}{3}\left(\frac{1}{R}+C_{R}\epsilon ^{p-2}\right).
	\end{split}\end{equation*}
Then, there exists a constant $\epsilon'(R)>0$ small enough such that when   $0<\epsilon\leq\epsilon'(R)$, one has
$$\frac{1}{R}+C_{R}\epsilon ^{p-2}\leq1 $$
and then
\begin{equation*}
	c+o_{n}(1)\geq \frac{1}{6}\|u_n\|^{2}_{H^{1}_{V}(\mathbb R^3)}- \frac{1}{6}\int_{\mathbb{R}^3}m u^{2}_n\text{d}x.
	\end{equation*}
Multiplying both sides by $\rho_{n}^{-2}$, for large $n$, we have 	
	$$\|v_n\|_{L^{2}(\mathbb{R}^3)}^{2}\geq \frac{1}{2m}.$$
Note that $v_n\to v$ strongly in $L^{2}(\mathbb{R}^3)$. Then $v\neq0$. Thus the Lebesgue measure of the set
$$\mathcal{M}=\{x\in\mathbb{R}^3:v(x)\neq 0\}$$
is positive. For $x\in\mathcal{M}$, as $n\to\infty$, we deduce
$$v_n(x)\to v(x)\neq 0,~u_{n}(x)=v_{n}(x)\rho_n\to\infty~\text{and}~\phi_{u_n}(x)\to\infty.$$

On the other hand, it follows from assumptions $({\bf f})$ and $({\bf f'})$ that for $\epsilon\in (0,\frac{\delta_0}{R})$,
\begin{equation}\label{eq21}
   F\left (\epsilon \eta _{R} \left ( \left | t\right |\right )t\right )\geq 0\ ~\text{for all}~t\in\mathbb{R}.
\end{equation}
From \eqref{eq21} and the Fatou lemma, for $\epsilon\in (0,\frac{\delta_0}{R})$ we have
\begin{equation*}
	\begin{split}
&\indent\int_{\mathbb{R}^3}\left(\frac{\phi_{u_n}u^{2}_{n}}{4}+\frac{1}{\epsilon^2}\eta _{R}\left (\|u_n\|^{2}_{H^1_{V}(\mathbb{R}^3)}\right )F\left (\epsilon \eta _{R} \left ( \left | u_n\right |\right )u_n\right )+\frac{1}{2}mu_{n}^{2}\right)\rho^{-2}_n\text{d}x\\
&\geq\int_{\mathbb{R}^3}\frac{\phi_{u_n}u^{2}_{n}}{4u_{n}^{2}}v^{2}_{n}\text{d}x\\
&\geq \frac{1}{4}\int_{\mathcal{M}}\phi_{u_n}v^{2}_n\text{d}x\to\infty,~~\text{as}~n\to\infty.
\end{split}
\end{equation*}
Thus
\begin{equation*}
		\begin{split}
		c-1&\leq I_{\epsilon,R}(u_n)\\
		&=\frac{1}{2}\int_{\mathbb{R}^3}\left(|\nabla u_n|^{2}+\tilde V(x)u_n^{2}\right)\text{d}x-\frac{1}{4}\int_{\mathbb{R}^3}\phi_{u_n}u_n^{2}\text{d}x\\
		&\indent-\frac{1}{\epsilon^{2}}\eta _{R}\left (\|u_n\|^{2}_{H^1_{V}(\mathbb{R}^3)}\right )\int_{\mathbb{R}^3}F\left (\epsilon \eta _{R} \left ( \left | u_n\right |\right )u_n\right )\text{d}x-\frac{m}{2}\int_{\mathbb{R}^3}u_n^{2}\text{d}x\\
		&=\rho_{n}^{2}\left(\frac{1}{2}-\int_{\mathbb{R}^3}\left(\frac{\phi_{u_n}u^{2}_{n}}{4}+\frac{1}{\epsilon^2}\eta _{R}\left (\|u_n\|^{2}_{H^1_{V}(\mathbb{R}^3)}\right )F\left (\epsilon \eta _{R} \left ( \left | u_n\right |\right )u_n\right )+\frac{1}{2}mu_{n}^{2}\right)\rho^{-2}_n\text{d}x\right)\\
        &\to-\infty,~~\text{as}~n\to\infty.
	\end{split}
    \end{equation*}
This is impossible. Therefore, our claim is true.

From the boundedness of $\{u_n\}$ and $H^{1}_{V}(\mathbb{R}^{3})$ is compactly embedded in $L^q(\mathbb{R}^3)(q\in[2,6))$, up to a sub-sequence we may assume
\begin{equation*}
		\begin{split}
	u_n&\to u_0~\text{weakly in}~H^{1}_{V}(\mathbb{R}^3);\\
	u_n&\to u_0~\text{strongly in}~L^{q}(\mathbb{R}^3)(q\in[2,6));\\
	u_n&\to u_0~\text{a.e. in}~ \mathbb{R}^3.
	\end{split}\end{equation*}
Next we prove the strong convergence of $u_{n}$ in $H_{V}^{1}\left ( \mathbb{R}^{3} \right ) $. To simplify the notation, let $$ \Psi_{\epsilon,R}\left ( u \right ):=\frac{1}{\epsilon ^{2}}\eta _{R}\left (\|u\|_{H^1_V(\mathbb{R}^3)}^{2}\right )\int_{\mathbb{R}^{3}}F\left (\epsilon \eta _{R} \left ( \left | u\right |\right )u\right )\text{d}x.$$

{\bf Claim}:
\begin{equation*}
    \begin{aligned}
     &\langle I'_{\epsilon,R}(u_n) - I'_{\epsilon,R}(u_0), u_n - u_0 \rangle \\
&= \|u_n - u_0\|_{H^1_V(\mathbb{R}^3)}^2-\int_{\mathbb{R}^3} (\phi_{u_n}u_n - \phi_{u_0}u_0)(u_n - u_0) \text{d}x -m\int_{\mathbb{R}^3}(u_n-u_0)^2\text{d}x\\
&\quad- \langle \Psi'_{\epsilon,R}(u_n) - \Psi'_{\epsilon,R}(u_0), u_n - u_0 \rangle\\
&\geq \frac{1}{2}\|u_n-u_0\|_{H^{1}_V(\mathbb{R}^3)}^{2}+o_n(1).
    \end{aligned}
\end{equation*}
Indeed, by the boundedness of ${u_n}$ in $H^{1}_{V}(\mathbb{R}^3)$, we have
\begin{equation*}
    \left |\int_{\mathbb{R}^{3}}\left ( \phi _{u_{n}} u_{n}-\phi _{u_{0}}u_{0}\right )\left ( u_{n}-u_{0} \right ) \text{d}x\right |\leq o_{n}\left ( 1 \right )\left\| u_{n}-u_{0}\right\|_{H^{1}_{V}(\mathbb{R}^3)}=o_{n}\left ( 1 \right )
\end{equation*}
and
\begin{equation*}
    \int_{\mathbb{R}^3}(u_n-u_0)^2\text{d}x=o_{n}(1).
\end{equation*}
Thus, we only need to prove that
\begin{equation*}
    \langle \Psi_{\epsilon,R}'(u_n) - \Psi_{\epsilon,R}'(u_0), u_n - u_0 \rangle \leq\frac{1}{2}\|u_n-u_0\|_{H^{1}_V(\mathbb{R}^3)}^{2}+o_n(1).
\end{equation*}
A direct computation shows that
\begin{align*}
\langle \Psi'_{\epsilon,R}(u), w \rangle =
& \frac{2}{\epsilon^2} \eta_R'(\|u\|_{H^1_{V}(\mathbb{R}^3)}^2) \int_{\mathbb{R}^3}(\nabla u\nabla w+\tilde V(x)uw) \text{d}x \int_{\mathbb{R}^3} F(\epsilon \eta_R(|u|) u)  \text{d}x \\
&+ \frac{1}{\epsilon} \eta_R(\|u\|_{H^1_{V}(\mathbb{R}^3)}^2) \int_{\mathbb{R}^3} f(\epsilon \eta_R(|u|) u) \left(\eta_R(|u|) + \eta_R'(|u|)|u| \right) w \text{d}x.
\end{align*}
And then
\begin{align}
      &\left<  \Psi'_{\epsilon,R} \left ( u_{n} \right )-\Psi_{\epsilon ,R}'\left ( u_{0} \right ),u_{n}-u_{0}\right>\nonumber \\
      &= \frac{2}{\epsilon^2} \eta_R'(\|u_n\|_{H^1_{V}(\mathbb{R}^3)}^2) \int_{\mathbb{R}^3}(\nabla u_n\nabla (u_{n}-u_{0})+\tilde V(x)u_n(u_{n}-u_{0})) \text{d}x \int_{\mathbb{R}^3} F(\epsilon \eta_R(|u_n|) u_n) \text{d}x \tag{T1}\\
     &\quad~- \frac{2}{\epsilon^2} \eta_R'(\|u_0\|_{H^1_{V}(\mathbb{R}^3)}^2) \int_{\mathbb{R}^3}(\nabla u_0\nabla (u_{n}-u_{0})+\tilde V(x)u_0(u_{n}-u_{0})) \text{d}x \int_{\mathbb{R}^3} F(\epsilon \eta_R(|u_0|) u_0)  \text{d}x \tag{T2} \\
    &\quad~+ \frac{1}{\epsilon} \eta_R(\|u_n\|_{H^1_{V}(\mathbb{R}^3)}^2) \int_{\mathbb{R}^3} f(\epsilon\eta_R(|u_n|) u_n) \left(\eta_R(|u_n|) + \eta_R'(|u_n|)|u_n| \right) (u_n-u_0) \text{d}x  \tag{T3}\\
    &\quad~- \frac{1}{\epsilon} \eta_R(\|u_0\|_{H^1_{V}(\mathbb{R}^3)}^2) \int_{\mathbb{R}^3} f(\epsilon \eta_R(|u_0|) u_0) \left(\eta_R(|u_0|) + \eta_R'(|u_0|)|u_0| \right) (u_n-u_0) \text{d}x. \tag{T4}
\end{align}
We now proceed to provide term-by-term estimates for $(\text{T}1)$ through $(\text{T}4)$. We begin by estimating $(\text{T}1)$, for which we obtain the following estimate:
\begin{equation*}
    \begin{aligned}
       \left | (\text{T}1)\right |&=\frac{2}{\epsilon^2} \eta_R'(\|u_n\|_{H^1_{V}(\mathbb{R}^3)}^2) \int_{\mathbb{R}^3}(\nabla u_n\nabla (u_{n}-u_{0})+\tilde V(x)u_n(u_{n}-u_{0})) \text{d}x \int_{\mathbb{R}^3} F(\epsilon \eta_R(|u_n|) u_n) \text{d}x\\
       &=\frac{2}{\epsilon^2} \eta_R'(\|u_n\|_{H^1_{V}(\mathbb{R}^3)}^2) \int_{\mathbb{R}^3}|\nabla (u_n-u_0)|^2+\tilde V(x)(u_{n}-u_{0})^2) \text{d}x \int_{\mathbb{R}^3} F(\epsilon \eta_R(|u_n|) u_n) \text{d}x\\
       &\quad+\frac{2}{\epsilon^2} \eta_R'(\|u_n\|_{H^1_{V}(\mathbb{R}^3)}^2) \int_{\mathbb{R}^3}(\nabla u_0\nabla (u_{n}-u_{0})+\tilde V(x)u_0(u_{n}-u_{0})) \text{d}x \int_{\mathbb{R}^3} F(\epsilon \eta_R(|u_n|) u_n) \text{d}x\\
       &\leq  \frac{2}{\epsilon^2} \eta_R'(\|u_n\|_{H^1_{V}(\mathbb{R}^3)}^2)\|u_n-u_0\|^2_{H^{1}_{V}(\mathbb{R}^3)}\int_{\mathbb{R}^3} F(\epsilon \eta_R(|u_n|) u_n)  \text{d}x\\
       &\quad+o_n(1)\frac{2}{\epsilon^2} \eta_R'(\|u_n\|_{H^1_{V}(\mathbb{R}^3)}^2)\int_{\mathbb{R}^3} F(\epsilon \eta_R(|u_n|) u_n) \text{d}x\\
      & \leq \frac{2}{\epsilon^2} \eta_R'(\|u_n\|_{H^1_{V}(\mathbb{R}^3)}^2) \|u_n-u_0\|^2_{H^{1}_{V}(\mathbb{R}^3)}\int_{\mathbb{R}^{3}}\left ( \delta \epsilon ^{2} \eta _{R}^{2}\left ( \left | u_n\right | \right )u_n^{2}+C_{\delta,R}\epsilon ^{p}\eta _{R}^{p}\left ( \left | u_n\right | \right )\left | u_n\right |^{p}\right )\text{d}x\\
      &\quad+o_{n}(1)\frac{2}{\epsilon^2} \eta_R'(\|u_n\|_{H^1_{V}(\mathbb{R}^3)}^2)\int_{\mathbb{R}^{3}}\left ( \delta \epsilon ^{2} \eta _{R}^{2}\left ( \left | u_n\right | \right )u_n^{2}+C_{\delta,R}\epsilon ^{p}\eta _{R}^{p}\left ( \left | u_n\right | \right )\left | u_n\right |^{p}\right )\text{d}x\\
      &\leq\delta \eta_R'(\|u_n\|_{H^1_{V}(\mathbb{R}^3)}^2) \|u_n-u_0\|^2_{H^{1}_{V}(\mathbb{R}^3)}\int_{\mathbb{R}^{3}}u_n^{2}\text{d}x \\
      &\quad+C_{\delta,R}\epsilon ^{p-2}\eta_R'(\|u_n\|_{H^1_{V}(\mathbb{R}^3)}^2)\|u_n-u_0\|^2_{H^{1}_{V}(\mathbb{R}^3)}\int_{\mathbb{R}^{3}}\left | u_n\right |^{p}\text{d}x\\
      &\quad+(\delta \eta_R'(\|u_n\|_{H^1_{V}(\mathbb{R}^3)}^2)\int_{\mathbb{R}^{3}}u_n^{2}\text{d}x+C_{\delta,R}\epsilon ^{p-2}\eta_R'(\|u_n\|_{H^1_{V}(\mathbb{R}^3)}^2)\int_{\mathbb{R}^{3}}\left | u_n\right |^{p}\text{d}x)o_{n}(1) \\
    &\leq \delta R^2\|u_n-u_0\|_{H^{1}_{V}(\mathbb{R}^3)}^{2}+C_{\delta,R}\epsilon ^{p-2}\|u_n-u_0\|_{H^{1}_{V}(\mathbb{R}^3)}^{2}+(\delta R^2+C_{\delta,R}\epsilon ^{p-2})o_{n}(1),
    \end{aligned}
\end{equation*}
where we use the fact that $u_n\to u_0$ weakly in $H^1_{V}(\mathbb{R}^3)$. Take $\epsilon(R)=\min\{\epsilon'(R),\frac{\delta_0}{R},\left(\frac{1}{4C_{\delta,R}}\right)^{1/(p-2)}\}$ with $\delta=\frac{1}{4R^2}$, for any $\epsilon\in(0,\epsilon(R))$ we can deduce that $$|(\text{T}1)|\leq \frac{1}{2}\|u_n-u_0\|_{H^{1}_{V}(\mathbb{R}^3)}^{2}+o_{n}(1).$$
Next, by using $u_{n} \to u _{0}~ \text{weakly~in}~H_{V}^{1}\left ( \mathbb{R}^{3}\right )$ and the similar method in Lemma \ref{lem3.5}, we obtain that
$$|(\text{T}2)|=o_{n}\left ( 1 \right ).$$

We now proceed to $(\text{T}3)$ and $(\text{T}4)$. Without loss of generality, we assume that $\epsilon(R)\in(0,1)$, for any $\epsilon\in(0,\epsilon(R))$, one has
\begin{equation*}
    \begin{aligned}
         \left |(\text{T}3)\right |&\leq   \frac{C_{R}}{\epsilon} \eta_R(\|u_n\|_{H^1_V(\mathbb{R}^3)}^2) \left|\int_{\mathbb{R}^3} f(\epsilon \eta_R(|u_n|) u_n) (u_n - u_0) \text{d}x\right|\\
    &\leq   C_{R}\eta_R(\|u_n\|_{H^1_V(\mathbb{R}^3)}^2) \int_{\mathbb{R}^3} \left [ \delta \eta _{R}\left ( \left | u_{n}\right | \right )|u_{n}|+C_{\delta,R } \epsilon^{p-2}\eta _{R}^{p-1}\left ( \left | u_{n}\right |\right )|u_{n}| ^{p-1}\right ]|u_n - u_0|\text{d}x\\
    &\leq\delta C_{R}\eta_R(\|u_n\|_{H^1_V(\mathbb{R}^3)}^2)   \int_{\mathbb{R}^3}|u_{n}||u_{n}-u_{0} |\text{d}x\\
    &\quad+ C_{\delta ,R}\epsilon^{p-2}\eta_R(\|u_n\|_{H^1_V(\mathbb{R}^3)}^2) \int_{\mathbb{R}^3}|u_{n}|^{p-1}|u_{n}-u_{0}|\text{d}x.\\
    &\leq \delta C_{R}\eta_R(\|u_n\|_{H^1_V(\mathbb{R}^3)}^2) \|u_n\|_{L^{2}(\mathbb{R}^3)}\|u_n-u_0\|_{L^{2}(\mathbb{R}^3)}\\
    &\quad+C_{\delta ,R}\epsilon^{p-2}\eta_R(\|u_n\|_{H^1_V(\mathbb{R}^3)}^2) \left ( \int_{\mathbb{R}^3}|u_{n}|^{p} \text{d}x\right )^{\frac{p-1}{p}}\left ( \int_{\mathbb{R}^3}\left | u_{n}-u_{0}\right |^{p}\text{d}x \right )^{\frac{1}{p}}\\
    &\leq \delta C_{R}\|u_n-u_0\|_{L^{2}(\mathbb{R}^3)}+C_{\delta ,R}\|u_n-u_0\|_{L^{p}(\mathbb{R}^3)}.
    \end{aligned}
\end{equation*}
Then it follows from $u_n\to u_0$ strongly in $L^{q}(\mathbb{R}^3)$ with $q\in[2,6)$ that $|(\text{T}3)|=o_{n}\left ( 1 \right ).$ Finally, it is easy to see that $|(\text{T}4)|=o_{n}\left ( 1 \right ).$

Combining these results, we have proved the above claim, and then one has
\begin{equation*}
    \|u_n - u_0\|_{H^1_V(\mathbb{R}^3)}^2=o_{n}\left ( 1 \right ),
\end{equation*}
which implies $u_{n} \to u _{0}$ strongly in $H_{V}^{1}\left ( \mathbb{R}^{3} \right )$.
\end{proof}

\end{document}
